\section{Introduction}
A self-avoiding path at its critical weight can make a long excursion
without paying an exponential cost. A basic quantitative question is how
much of its weight reaches the opposite side of a strip. On the regular
honeycomb lattice we prove that this crossing mass has power $-1/4$,
with upper and lower bounds by fixed positive constants. Our proof also
determines the first displacement moment of paths that return to the
starting side. This second observable provides the route to the crossing
estimate.

\subsection{The model and the result}
We use the honeycomb lattice dual to the equilateral triangular tiling
of side length one, with one triangular edge direction horizontal.
The separation of consecutive horizontal band boundaries is
$d_0=\sqrt3/2$. A \emph{port} is the midpoint of an edge on the boundary
of a union of triangles. A path between distinct ports follows dual edges,
including the two terminal half-edges, visits each dual vertex at most
once, and meets the boundary only at its endpoints. Its length
$|\gamma|$ is the number of visited dual vertices; ports carry no weight.
We assign the critical weight $\rho^{|\gamma|}$, where
\begin{equation}\label{eq:critical-weight}
 \rho=\frac1{\sqrt{2+\sqrt2}},\qquad c=\cos(3\pi/8).
\end{equation}
This vertex convention agrees with the usual edge convention for the
connective constant up to the fixed endpoint factor. Duminil-Copin and
Smirnov proved that the reciprocal of $\rho$ is the honeycomb connective
constant~\cite{DCS2012}.

Let $\mathcal S_N$ be the infinite horizontal strip of $N$ triangular
bands, where $N\ge1$. Index its bottom ports consecutively by $\mathbb Z$,
fix the source at port $0$, and write $Z_D(a,b)$ for the sum of critical
weights of paths from $a$ to $b$ in $D$. A path returning to the bottom
is called an \emph{arch}; a path ending at the top is called a
\emph{bridge}. In both cases all other points of the path are in the
interior of the strip. Define
\begin{equation}\label{eq:01-strip-quantities}
 \begin{split}
 K_N(k)&=Z_{\mathcal S_N}(0,k),\\
 \mathcal A_N&=\sum_{k\ne0}K_N(k),\qquad
 \mathcal B_N=\sum_{b\text{ on the top}}Z_{\mathcal S_N}(0,b),\\
 m_N&=\sum_{k\ge1}kK_N(k).
 \end{split}
\end{equation}
The sums run over actual admissible boundary ports. Adjacent bottom
ports have horizontal separation one, so $m_N$ is the first rightward
horizontal-displacement moment. Top ports lie on the translated lattice
$N/2+\mathbb Z$. These are unnormalized masses: for example the bridge
probability measure would divide each bridge weight by $\mathcal B_N$.
Figure~\ref{fig:model-conventions} shows the convention.

\begin{figure}[htbp]
\centering
\begin{tikzpicture}[x=1.55cm,y=1.55cm,font=\small,>=Latex]
 \pgfmathsetmacro{\bandheight}{sqrt(3)/2}
 \fill[blue!4] (-.8,\bandheight) rectangle (3.2,2*\bandheight);
 \begin{scope}
  \clip (-.8,0) rectangle (3.2,3*\bandheight);
  \foreach \j in {0,1,2,3}{
   \draw[black!22] (-1,\j*\bandheight)--(4,\j*\bandheight);
  }
  \foreach \j in {0,1,2}{
   \foreach \i in {-3,-2,...,4}{
    \draw[black!22] ({\i+\j/2},\j*\bandheight)
      --({\i+\j/2+.5},{(\j+1)*\bandheight});
    \draw[black!22] ({\i+\j/2},\j*\bandheight)
      --({\i+\j/2-.5},{(\j+1)*\bandheight});
   }
  }
 \end{scope}
 \draw[thick] (-.8,0)--(3.2,0);
 \draw[thick] (-.8,3*\bandheight)--(3.2,3*\bandheight);
 \coordinate (a) at (.5,0);
 \coordinate (v1) at (.5,\bandheight/3);
 \coordinate (v2) at (1,2*\bandheight/3);
 \coordinate (v3) at (1,4*\bandheight/3);
 \coordinate (v4) at (.5,5*\bandheight/3);
 \coordinate (v5) at (.5,7*\bandheight/3);
 \coordinate (v6) at (1,8*\bandheight/3);
 \coordinate (b) at (1,3*\bandheight);
 \draw[very thick,blue!65!black] (a)--(v1)--(v2)--(v3)--(v4)--(v5)--(v6)--(b);
 \foreach \v in {v1,v2,v3,v4,v5,v6}{\fill (\v) circle (1.4pt);}
 \foreach \v in {a,b}{\draw[thick,fill=white] (\v) circle (2pt);}
 \node[below=3pt] at (a) {initial port};
 \node[above=3pt] at (b) {terminal port};
 \node[align=left,anchor=west] (vertex) at (1.65,1.05) {visited dual\\vertex};
 \draw[->] (vertex.west)--(v3);
 \draw[<->] (3.4,\bandheight)--node[right]{one band}(3.4,2*\bandheight);
 \draw[black!50] (3.2,\bandheight)--(3.5,\bandheight);
 \draw[black!50] (3.2,2*\bandheight)--(3.5,2*\bandheight);
\end{tikzpicture}
\caption{Port and length conventions. The path passes through six adjacent
triangles of the triangular tiling, so it visits six vertices of the
honeycomb dual and has port weight $\rho^6$. Its two terminal pieces are
half-edges. An ordinary vertex-to-vertex path counts full edges instead. The shaded region is one triangular band.}
\label{fig:model-conventions}
\end{figure}



\begin{theorem}[Critical strip mass]\label{thm:strip}
For every integer $N\ge1$, all sums in
\eqref{eq:01-strip-quantities} are finite. With constants independent of
$N$,
\begin{equation}\label{eq:01-strip-result}
 c\mathcal A_N+\mathcal B_N=1,\qquad
 m_{N+1}-m_N\asymp\mathcal B_N,\qquad
 m_N\asymp N^{3/4},\qquad
 \mathcal B_N\asymp N^{-1/4}.
\end{equation}
Moreover, $\mathcal B_N$ is nonincreasing in $N$.
\end{theorem}
Here $f_N\asymp g_N$ means that their ratio lies between two fixed
positive constants. Thus the theorem determines a bounded-factor
power law, rather than only a logarithmic exponent. If a height-zero
term is useful, we set $\mathcal B_0=1$ by convention; no strip path is
being added to the definition above. The crossing estimate can then be
written $\mathcal B_N\asymp(1+N)^{-1/4}$ for $N\ge0$.

\subsection{Earlier work}
The two-dimensional polymer exponents arose from the dilute $O(n)$
model and its Coulomb-gas analysis. Nienhuis predicted the honeycomb
critical point and the self-avoiding-walk exponents at
$n=0$~\cite{Nienhuis1982}.
Lawler, Schramm and Werner developed the conjectural conformally
invariant scaling picture for planar self-avoiding walk
\cite[Sections~3.3.1, 3.4.3 and~4.1]{LawlerSchrammWerner2004}.
Their boundary exponent $5/8$ predicts a boundary-to-boundary mass
of power $-5/4$. Summing the endpoint along the opposite side of a
strip then predicts the crossing power $-1/4$; this deduction is
also stated explicitly in \cite[Section~4]{DCS2012}.

Duminil-Copin and Smirnov proved the critical-point prediction using
a parafermionic observable whose local cancellation gives a positive
boundary identity~\cite[Lemmas~1 and~2]{DCS2012}.
Their strip argument gives bounds of order $1/N$ and $1$ for the
crossing mass. These bounds leave the predicted power undetermined.

The decay of this mass was subsequently proved by Beaton,
Bousquet-M\'elou, de Gier, Duminil-Copin and Guttmann in their study of
surface adsorption~\cite[Theorem~10]{BeatonEtAl2014}.
Glazman and Manolescu obtained a shorter proof, a logarithmic bound
along a sequence of heights, and invariance of the half-plane boundary
two-point function under columnwise rhombic deformations with angles
in $[\pi/3,2\pi/3]$
\cite[Proposition~1.1 and Theorems~1--2]{GM2019}.
Krachun and Panagiotis proved a polynomial upper bound
$\mathcal B_N\le100N^{-10^{-10}}$ in their strip convention,
together with quantitative sub-ballisticity
\cite[Theorems~1--2]{KrachunPanagiotis2026}.
Theorem~\ref{thm:strip} establishes the predicted strip power with
uniform multiplicative constants in the port convention above.
The exponent $3/4$ for $m_N$ describes an unnormalized displacement
moment summed over all path lengths, with strip height as its scale.
The fixed-length displacement conjecture and the conjectured conformal
scaling limit concern different observables.

\subsection{Proof strategy}
The boundary identity alone does not determine this power. Our route
is to calculate $m_N$ first and then recover $\mathcal B_N$ from its
increments. The argument passes from finite strip transfers to a
polynomial formula, then to a positive integral that can be estimated
uniformly in the height.

In Section~\ref{sec:flux}, cancellation of closed excursions gives the
boundary identity and controls finite-height transfer matrices. We
introduce rhombic weights with one spectral parameter in each row.
A cut between columns records vacant positions and occupied positions paired
by path pieces on one side of the cut. The stationary vacuum vector
records the weighted sums of these partial pairings with no exterior
source; a source vector has one additional strand connected to a
specified boundary source.
Their local relations belong to the dilute-loop integrable structure;
the relation with discrete parafermions was developed by Ikhlef and
Cardy, and the weighted-walk normalization and corrected general
formulas were treated by Glazman~\cite{IkhlefCardy2009,Glazman2015}.
We prove the exact local identities used here in
Appendix~\ref{app:local-calculations}.

Section~\ref{sec:vacuum} constructs a scalar Pfaffian $p_N$ that turns
the stationary vacuum into a vector of bounded-degree polynomials.
A degree bound is essential: it makes interpolation an identity
argument. Exchange relations and reductions at special parameters
have this role in earlier dense and dilute loop calculations
\cite{DiFrancescoZinnJustin2005,GarbaliNienhuis2017Ground,GarbaliNienhuis2017Sum}.
Here the degree and stationarity arguments are proved for zero loop
weight.
Section~\ref{sec:arch} cuts an arch at every column between its
endpoints. This converts $m_N$ into a pairing of source vectors and
then into one coefficient of $p_{N+1}/p_N$. A stepped-boundary flux
identity proves $m_{N+1}-m_N\asymp\mathcal B_N$.

Section~\ref{sec:integral} converts the homogeneous polynomial problem
into a positive integral over $0<y_1<\cdots<y_N<1$, with density
proportional to
\[
 \prod_{i=1}^N\omega(y_i)
 \prod_{i<j}\frac{(y_j-y_i)^2}{y_i+y_j},\qquad
 \omega(y)=\frac12\sqrt{\sqrt{\frac{1+y}{2y}}-1}.
\]
Positivity of a Cauchy-moment determinant makes the coincident-parameter
limit nonsingular.
The pair interaction in the resulting integral is the one appearing in
generalized Bures ensembles~\cite{ForresterKieburg2016}.
The derivation from the strip polynomial fixes the normalization and
identifies the integral with the arch moment. Finally,
Section~\ref{sec:hard-edge} compares this measure with Bures--Laguerre
laws. Schur's Pfaffian identity and de Bruijn's integration formula
give their normalizing constants; positive association gives the
comparison. The smallest-coordinate estimates yield
$m_N\asymp N^{3/4}$. Monotonicity and the increment identity then give
the crossing power in Theorem~\ref{thm:strip}.
